\chapter{chapitre 1} 
\begingroup
\parindent=0em
\etocsettocstyle{\rule{\linewidth}{\tocrulewidth}\vskip0.5\baselineskip}{\rule{\linewidth}{\tocrulewidth}}
\localtableofcontents 
\endgroup
\newpage



\section{Introduction}
\section{Tableau}
\begin{table}[!h]
    \renewcommand{\arraystretch}{1.2}
    \caption{Variables définies.}
\centering
\begin{tabular}{|l|p{10.5cm}|}
        \hline
        \rowcolor[HTML]{ECF4FF} 
        \textbf{Variables Déclarées} & \textbf{Description} \\ \hline
        \textbf{network[de][à]} & Une file d'attente de messages envoyés d'un processus à un autre. \\ \hline
        \textbf{activeClients} & Une liste contenant les identifiants des clients actifs. \\ \hline
        \textbf{Topics} & Une liste de sujets disponibles. \\ \hline
        \textbf{TopicPool} & Une liste associant chaque sujet à un ensemble de clients qui y sont abonnés. \\ \hline
        \textbf{{SessionClients}} & Une liste d'identifiants des clients avec une session active. \\ \hline
        \textbf{{RetainedMessages}} & Une liste contenant les messages conservés par le courtier. \\ \hline
    \end{tabular}
    \label{variables}
\end{table}
\FloatBarrier

\section{Conclusion}
